%% =============================================================================
%% configuracoes.tex — pacotes e comandos adicionais (carregado pelo main.tex)
%% A classe utfpr.cls já carrega abntex2/memoir, babel, hyperref, graphicx, amsmath,
%% caption, glossaries, nomencl, tocloft, titlesec e ulem, entre outros.
%% =============================================================================

%% Pacotes adicionais
\usepackage{bigdelim, booktabs, colortbl, longtable, multirow}%% Tabelas
\usepackage{amssymb, amstext, amsthm, icomma}%% Matemática (icomma: vírgula decimal)
\usepackage{pifont, textcomp, wasysym}%% Símbolos de texto
\usepackage{pdfpages}%% Inserir páginas de PDF (\includepdf)
\usepackage{xspace}
\usepackage{xcolor}

%% Listagens de código (Python/Java) do trabalho
\usepackage{listings}
\lstset{
  language=Python,
  basicstyle=\ttfamily\footnotesize,
  keywordstyle=\color{blue!70!black}\bfseries,
  commentstyle=\color{green!45!black}\itshape,
  stringstyle=\color{red!60!black},
  numbers=left, numberstyle=\tiny\color{gray},
  frame=single, breaklines=true, showstringspaces=false,
  tabsize=4, captionpos=t,
  extendedchars=true, inputencoding=utf8,
  literate={á}{{\'a}}1 {ã}{{\~a}}1 {â}{{\^a}}1 {à}{{\`a}}1 {é}{{\'e}}1 {ê}{{\^e}}1
           {í}{{\'i}}1 {ó}{{\'o}}1 {õ}{{\~o}}1 {ô}{{\^o}}1 {ú}{{\'u}}1 {ç}{{\c{c}}}1
           {Á}{{\'A}}1 {É}{{\'E}}1 {Ç}{{\c{C}}}1
}
\renewcommand{\lstlistingname}{Código}

%% Marcador de pendência visível no PDF: \todo{texto}
\newcommand{\todo}[1]{\textcolor{red}{\textbf{[A FAZER:} #1\textbf{]}}}

%% Caixa de aviso para anotações do grupo: \caixa{Título}{texto}
\usepackage{tcolorbox}
\newcommand{\caixa}[2]{%
  \begin{tcolorbox}[colback=red!5!white, colframe=red!45!white, title=#1, fonttitle=\bfseries]
  #2
  \end{tcolorbox}%
}

%% Seção quaternária (\subsubsection) sublinhada no sumário e no título (padrão UTFPR)
\renewcommand*{\cftsubsubsectionfont}{\normalfont\uline}
\setsubsubsecheadstyle{\ABNTEXsubsubsectionfont\ABNTEXsubsubsectionfontsize\ABNTEXsubsubsectionupperifneeded\uline}

%% Comandos personalizados
\newcommand{\cpp}{\texttt{C$++$}}%% C++
\newcommand{\latex}{\LaTeX\xspace}%% LaTeX
\newcommand{\ds}{\displaystyle}%% Equação em tamanho normal
\newcommand{\bsym}[1]{\boldsymbol{#1}}%% Negrito em modo matemático
\newcommand{\mr}[1]{\mathrm{#1}}%% Texto reto em modo matemático
\newcommand{\der}{\mr{d}}%% Operador diferencial
\newcommand{\deri}[2]{\frac{\der #1}{\der #2}}%% Derivada ordinária
\newcommand{\derip}[2]{\frac{\partial #1}{\partial #2}}%% Derivada parcial
\newcommand{\pare}[1]{\left( #1 \right)}%% Parênteses
\newcommand{\colc}[1]{\left[ #1 \right]}%% Colchetes
\newcommand{\chav}[1]{\left\lbrace #1 \right\rbrace}%% Chaves
\newcommand{\sen}{\operatorname{sen}}%% Seno
\newcommand{\senh}{\operatorname{senh}}%% Seno hiperbólico
\newcommand{\tg}{\operatorname{tg}}%% Tangente
\newcommand{\tgh}{\operatorname{tgh}}%% Tangente hiperbólica
\newcommand{\seqref}[1]{Equação~\eqref{#1}}%% Referência a uma equação
\newcommand{\meqref}[1]{Equações~\eqref{#1}}%% Referência a várias equações
%% Atalhos de citação (abntex2cite): \citep = \cite, \citet = \citeonline
\newcommand{\citep}[1]{\cite{#1}}
\newcommand{\citet}[1]{\citeonline{#1}}
\newcommand{\citepa}[1]{(\citeauthor{#1})}%% Só autor, entre parênteses
\newcommand{\citeta}[1]{\citeauthoronline{#1}}%% Só autor, no texto
\newcommand{\citepy}[1]{(\citeyear{#1})}%% Só ano, entre parênteses
\newcommand{\citety}[1]{\citeyear{#1}}%% Só ano, no texto

%% Família da fonte do texto (usado em variaveis.tex: \fonteTexto{\sfdefault} ou {\rmdefault})
\newcommand{\fonteTexto}[1]{\renewcommand{\familydefault}{#1}}

%% Pasta das figuras
\graphicspath{{figuras/}}

%% Evita linhas órfãs e viúvas
\widowpenalty=10000
\clubpenalty=10000

%% Recuo das listas (itemize/enumerate) de 1,5 cm, igual ao parágrafo
\setlength{\leftmargini}{1.5cm}
\setlength{\leftmarginii}{1.5cm}

